\documentclass[a4paper,
fontsize=11pt,
%headings=small,
oneside,
numbers=noperiodatend,
parskip=half-,
bibliography=totoc,
final
]{scrartcl}

\usepackage[babel]{csquotes}
\usepackage{synttree}
\usepackage{graphicx}
\setkeys{Gin}{width=.4\textwidth} %default pics size

\graphicspath{{./plots/}}
\usepackage[ngerman]{babel}
\usepackage[T1]{fontenc}
%\usepackage{amsmath}
\usepackage[utf8x]{inputenc}
\usepackage [hyphens]{url}
\usepackage{booktabs} 
\usepackage[left=2.4cm,right=2.4cm,top=2.3cm,bottom=2cm,includeheadfoot]{geometry}
\usepackage[labelformat=empty]{caption} % option 'labelformat=empty]' to surpress adding "Abbildung 1:" or "Figure 1" before each caption / use parameter '\captionsetup{labelformat=empty}' instead to change this for just one caption
\usepackage{eurosym}
\usepackage{multirow}
\usepackage[ngerman]{varioref}
\setcapindent{1em}
\renewcommand{\labelitemi}{--}
\usepackage{paralist}
\usepackage{pdfpages}
\usepackage{lscape}
\usepackage{float}
\usepackage{acronym}
\usepackage{eurosym}
\usepackage{longtable,lscape}
\usepackage{mathpazo}
\usepackage[normalem]{ulem} %emphasize weiterhin kursiv
\usepackage[flushmargin,ragged]{footmisc} % left align footnote
\usepackage{ccicons} 
\setcapindent{0pt} % no indentation in captions
\usepackage{xurl} % Breaks URLs

%%%% fancy LIBREAS URL color 
\usepackage{xcolor}
\definecolor{libreas}{RGB}{112,0,0}

\usepackage{listings}

\urlstyle{same}  % don't use monospace font for urls

\usepackage[fleqn]{amsmath}

%adjust fontsize for part

\usepackage{sectsty}
\partfont{\large}

%Das BibTeX-Zeichen mit \BibTeX setzen:
\def\symbol#1{\char #1\relax}
\def\bsl{{\tt\symbol{'134}}}
\def\BibTeX{{\rm B\kern-.05em{\sc i\kern-.025em b}\kern-.08em
    T\kern-.1667em\lower.7ex\hbox{E}\kern-.125emX}}

\usepackage{fancyhdr}
\fancyhf{}
\pagestyle{fancyplain}
\fancyhead[R]{\thepage}

% make sure bookmarks are created eventough sections are not numbered!
% uncommend if sections are numbered (bookmarks created by default)
\makeatletter
\renewcommand\@seccntformat[1]{}
\makeatother

% typo setup
\clubpenalty = 10000
\widowpenalty = 10000
\displaywidowpenalty = 10000

\usepackage[colorlinks, linkcolor=black,citecolor=black, urlcolor=libreas,
breaklinks= true,bookmarks=true,bookmarksopen=true]{hyperref}
\usepackage{hyperxmp}
\usepackage{breakurl}

%meta
%meta

\fancyhead[L]{N. Fötsch, M. Hordijk, I. Smudja\\ %author
LIBREAS. Library Ideas, 48 (2026). % journal, issue, volume.
\href{https://doi.org/10.18452/37900color{black}https://doi.org/10.18452/37900}
{}} % doi
\fancyhead[R]{\thepage} %page number
\fancyfoot[L] {\ccLogo \ccAttribution\ \href{https://creativecommons.org/licenses/by/4.0/}{\color{black}Creative Commons BY 4.0}}  %licence
\fancyfoot[R] {ISSN: 1860-7950}

\title{\LARGE{On the front lines of access: librarians discuss book bans in the Netherlands and the United States}}% title
\author{Norma Fötsch, Marjolein Hordijk, Ivana Smudja} % author

\setcounter{page}{1}

\hypersetup{%
      pdftitle={On the front lines of access: librarians discuss book bans in the Netherlands and the United States},
      pdfauthor={Norma Fötsch, Marjolein Hordijk, Ivana Smudja},
      pdfcopyright={CC BY 4.0 International},
      pdfsubject={LIBREAS. Library Ideas, 48 (2026).},
      pdfkeywords={Book Bans, censorship, Netherlands, United States of America, library},
      pdflicenseurl={https://creativecommons.org/licenses/by/4.0/},
      pdfcontacturl={http://libreas.eu},
      pdfurl={https://doi.org/10.18452/37900},
      pdfdoi={10.18452/37900},
      pdflang={en},
      pdfmetalang={en}
     }



\date{}
\begin{document}

\maketitle
\thispagestyle{fancyplain} 

%abstracts
\begin{abstract}
\noindent
\textbf{Kurzfassung}: In der Universitätsbibliothek von Nimwegen führte
die unabhängige Filmemacherin Ivana Smudja ein Interview mit zwei
Verfechterinnen des Zugangs zu Informationen: Marjolein Hordijk und
Norma Fötsch. Gemeinsam beleuchteten sie die Realität verbotener Bücher
in den Niederlanden und den Vereinigten Staaten, erörterten die
Beweggründe für Zensur und reflektierten über die Bedeutung der
Lesefreiheit.

\begin{center}\rule{0.5\linewidth}{0.5pt}\end{center}

\noindent\textbf{Abstract}: At the university library in Nijmegen, Ivana Smudja,
an independent film\-maker, met with bibliophiles and
access-to-information advocates Marjolein Hordijk and Norma Fötsch.
Together, they explored the reality of banned books in the Netherlands
and the United States, unpacked the motivations behind censorship, and
reflected on the importance of the freedom to read.
\end{abstract}

%body
\vspace{0.5cm}
What follows is an edited transcript of this conversation, published
alongside a short
\href{https://www.voxweb.nl/nieuws/waarom-worden-boeken-verboden-en-kan-dit-ook-in-nederland-gebeuren}{audiovisual
summary} that was released earlier.

Transcript by Norma Fötsch and Marjolein Hordijk

Video and questions by Ivana Smudja

\textbf{1. What can we see on the table here, and what makes these books
special or controversial?}

\emph{Norma Fötsch (NF):} My colleagues regularly organise exhibitions
to raise awareness about specific issues. The table next to us is an
example of this; it focuses on banned books and shows textbooks on
banned books and examples of banned books.

Throughout history, books have been banned or access to them has been
restricted. The reasons for these can vary greatly and include
political, religious or social reasons, or the fact that they contain
sexual content. The list of controversial books is incredibly long.

\textbf{2. Could you show us an example of a banned book? Why is this
book banned?}

\emph{Marjolein Hordijk (MH):} Banned books? I could bring my whole
bookcase. \enquote{The Hate U Give} is one of the most frequently banned
books in American school libraries, and one of my favourite books. The
reasons given include concerns about coarse language, drug use, and
sexual scenes and references. It is also claimed that the book is
anti-police and promotes a particular social agenda.

The book is about a girl, Starr, who grows up in \enquote*{the hood} but
attends a white school outside the neighbourhood because her parents
think it's better for her. As a result, she feels out of place in both
groups. Until the moment comes when her best friend is killed by a white
police officer who thought the boy was pulling out a weapon. It was
actually a hairbrush. A Black Lives Matter movement against the police
emerges. As the key witness, Starr eventually decides to speak out and
she becomes the figurehead of these demonstrations.

So you could argue that all the issues mentioned earlier, such as coarse
language and sexual scenes, are being used as grounds for removing a
book from a school library. But it certainly raises eyebrows. After all,
isn't it far more likely that people simply don't want books that
empower black communities? Because this doesn't fit with the
conservative white American narrative?

\textbf{3. What does it mean when a book is
\textquotesingle banned\textquotesingle? Does it refer to a legal
prohibition, or to the removal of books from libraries, for example?}

\emph{NF:} When a book is banned, it is worth checking precisely what is
prohibited. Is printing or distribution prohibited, or is access to the
book restricted? For example, if the ban concerns printing the book, it
may be that someone else holds the copyright, meaning you are not
permitted to do so. If the issue is more about the content, however, it
could be considered censorship. Censorship involves government
supervision of spoken and written expressions that it deems
objectionable. Measures can range from the removal of publications and
blocking access to the internet to intimidation and arrests.

In the Netherlands, you do not need prior permission to publish
anything. There is also no such thing as an index of banned books.
However, the distribution of texts can be prohibited by the courts after
publication, and authors can be prosecuted for inflammatory texts, for
example.

In the United States, ongoing book bans often involve the removal of
books from school libraries. Access to these books is restricted. These
restrictions are often initiated by non-governmental organisations, such
as Moms for Liberty. Therefore, book bans are not just
government-imposed, but community-led too.

\textbf{4. What is your opinion on how we should deal with difficult or
controversial books? Should they always be made available?}

\emph{NF:} When it comes to dealing with controversial books,
conflicting interests come into play. Let me introduce you to two
examples.

The first concerns the right to free speech versus anti-discrimination
law. Freedom of speech is an important right that should be respected.
But should we accept that anyone can publish or republish books in which
a Nazi spreads hateful thoughts and calls for the killing of anyone who
does not conform to the norm? No, we should not accept it, and you
probably know the book I\textquotesingle m talking about: \enquote{Mein
Kampf} by Adolf Hitler.

Another example of conflicting interests is the need for access to the
original text versus adapting it to modern language. Some
children\textquotesingle s books use the n-word. Should we destroy the
original books and only make the adapted versions available? This is a
real example. In Denmark, for instance, there was a significant debate
surrounding the adaptation of Halfdan Rasmussen\textquotesingle s
children\textquotesingle s poems. Ultimately, eight of his poems were
not republished due to their racist content, but the original book
containing all the poems is still available.

These examples demonstrate that the issue of banning books is complex.
But don\textquotesingle t get me wrong, governments and communities
would be wise to avoid banning books. A true democracy can tolerate
different perspectives.

There are plenty of other ways to deal with questionable publications.
As a publisher, for instance, you could publish an annotated edition of
a controversial book. As a library, you could offer textbooks alongside
the original book to explain why it is harmful or discriminatory. As a
reader, you could choose to read books that highlight different
perspectives.

\textbf{5. Are there any banned books or books with restricted access in
this library? How do you make that decision?}

\emph{NF:} The Radboud University Library\textquotesingle s core task is
to provide essential books and materials for study and research. This
distinguishes it from a public library, which serves a broader public
purpose. Furthermore, our aim is not to preserve every book, since we do
not fulfil an archival function. Our current collection has been
carefully developed in close consultation with the faculties. Every new
acquisition is subject to careful assessment to determine its potential
contribution to study and/or research. The criteria used for this
assessment may vary from publication to publication and from faculty to
faculty.

\textbf{6. What is the current climate in the US and Europe with regard
to sensitive or controversial titles?}

\emph{MH:} We must not assume that book bans only happen in other
countries. Here in Europe, too, we are seeing a strong shift towards
right-wing conservatism. Here, too, there are books that are under
attack. These are, in particular, books relating to gender and
sexuality. When populism is on the rise, there is no need for critical
thinkers. Books give people a broader perspective and encourage them to
think for themselves.

\textbf{7. In the United States, an increasing number of books are being
removed from school libraries. Is this indeed the case, and what kind of
books are we talking about?}

\emph{NF:} Unfortunately, I can be brief about this. This picture is
accurate. In the US, books are being removed from school libraries.
Since 2021, PEN America has been monitoring which books are being
removed. PEN America is an organisation dedicated to literature and
human rights. According to their records, almost 10,000 titles are no
longer permitted in school libraries. Notably, the majority of these
books are about people of colour, such as \enquote{The Hate U Give},
which Marjolein mentioned, as well as books about queer and political
issues.

An example of a queer book is: \enquote{Gender Queer: A Memoir}. It was
written and illustrated by Maia Kobabe. The protagonist discovers what
it means to be nonbinary and that they are romantically attracted to
others, but not sexually attracted to anyone. Despite this, the book has
been banned for its sexual content. The book \enquote{Fun Home: A Family
Tragicomic} by the great Alison Bechdel has also been banned for its
sexual content. The book is about the author's coming of age as a
lesbian and deals with gender roles and dysfunctional family life, among
other things. One final example I would like to share with you is
\enquote{Maus}, a graphic novel by Art Spiegelman. It recounts the
experiences of a Polish Jew who survived the Holocaust. It took me a
while to find out why this book was banned. It turned out that the
depiction of naked mice representing Jews in a concentration camp was
the reason for the ban.

Books are often banned on the grounds of sexual content. However, this
argument is applied very selectively in the US. It is primarily employed
to prohibit books that address unwelcome topics, such as racism, gay and
lesbian sexuality, and critical historical reflections.

\textbf{8. Are any books banned in the Netherlands? If so, which ones
and why? And if not so, have there ever been plans to ban certain
titles?}

\emph{MH:} At present times, no books are banned. \enquote{Mein Kampf}
was banned for a very long time, but that was partly because the book
could be kept off the shelves due to copyright. When the copyright
expired, there was a debate about what to do with the book because of
its hateful content. Ultimately, an annotated edition was published. I
think it is important to continue having these kinds of discussions.

In the Netherlands, it is enshrined in law that libraries are autonomous
in choosing the composition of their collections. The government cannot
intervene in this. However, social or religious pressure can influence
this. For instance, there are many schools with a religious basis that
refuse to include books, often gender-related, in their collections.
Libraries and schools are constantly discussing this issue. They are
seeking a balance between their own ideals and the reading enjoyment of
many children.

Although in practice, titles can only be banned if they are extremely
hateful and discriminatory, it is very rare for a book to actually be
banned. However, retailers such as BOL.com and NBD Biblion, a library
supplier, do make a selection. They do not actively offer those books.

\textbf{9. What are the implications of certain books being banned from
schools and universities? What consequences may this have?}

\emph{MH:} Put very simply, you can't be something you don't know about.
If you're a young boy and you feel that you're more attracted to other
boys, but you've never heard of \enquote*{being gay}, then you can't be
gay. Because you don't know what it is.

Banning books creates a very one-sided view of the world. But that
doesn't mean all these other feelings don't exist. It leads to one-sided
representation and unhappy children who don't fit the mould. That's a
lot of children. Anyone and everyone who doesn't fit the heteronormative
white mould won't recognise themselves. It's precisely those books that
provide mirrors and windows. In a mirror, you can recognise yourself,
and a window gives you a different view of the world.

When you talk about banning books by critical thinkers or influential
writers, you are erasing parts of history. You are erasing people's
identities. As far as I'm concerned, it's the first step towards
dehumanisation.

\textbf{10. Why is it so important for a democratic society to have free
access to books and knowledge?}

\emph{NF:} For a democratic society to thrive, people must understand
political decisions, question them critically, and take action
themselves. Books, journals and other academic and non-acade\-mic
publications provide background information and offer different
perspectives. Without this information, we cannot assess current events
critically, nor can we take action.

Nearly 4,000 unique titles are on last year\textquotesingle s list of
banned books. Most of the affected titles involve publications on
LGBTQIA+ topics, race, and politics. The aim is to prevent young people
from reading these books. The aim is to make these books invisible.
However, it\textquotesingle s not just about making these books
invisible; it\textquotesingle s also about making queer people, people
of colour, and left-wing individuals invisible.

On a personal level, I believe that free access to books is a necessity,
not a luxury. It is a prerequisite for a society in which diversity and
equal treatment are valued. For instance, removing queer books from
school libraries excludes people and renders queer individuals
invisible. As a child in the 1980s, I would have loved to read books
about two girls falling in love, for example. However, I couldn't find
any. Consequently, as an adult, I have struggled to find the right
words. Fortunately, these books now exist, and I am happy for young
people, hoping that these books will remain available to them.

\textbf{11. What is your view on the future? Do you think that more
books will disappear, or that banned titles will be reinstated?}

\emph{MH:} We haven't hit rock bottom yet. In the Netherlands, we're
seeing coordinated campaigns gaining traction, such as the
\enquote*{Week van de Lentekriebels} project for primary schools, which
introduces pupils to sex and relationship education, and campaigns
targeting certain authors. They're having the desired effect. People are
complaining to their children's schools and speaking out on social
media. There are political parties currently sitting in the House of
Representatives whose manifestos call for the desexualisation of school
textbooks. Many authors have now started to self-censor out of fear.
Publishers are also asking their authors to delete passages. We live in
a time when a children's book author received so many death threats that
he had to be put under security protection and leave his home. This is
already happening, and although I hope we will get past this very soon,
I fear it will take some time and is likely to get even worse.

\textbf{12. You are organising and participating in a screening of the
documentary \enquote{The Librarians} at the city film house, followed by
a discussion on book censorship in the US. Why is this discussion
important, and what do you hope to achieve with the event?}

\emph{NF:} I would like to use this opportunity to emphasise the
importance of libraries and free access to information. My personal aim
is to raise awareness of this issue. I want my colleagues and friends to
see what is happening, and the documentary
\enquote{\href{https://thelibrariansfilm.com/}{The Librarians}} by Kim
A. Snyder shows this in a really compelling way.

I have previously watched this documentary. Afterwards, I felt a mixture
of powerlessness and anger, but also a desire to take action. I am not
the only one who felt this way. Marjolein just told me that becoming
active is also the topic of the manifesto \enquote{Doe Iets}, published
last year by ProBiblio. It\textquotesingle s about activism in
libraries. When I think about what I can do, a lot of activities come to
mind. You could think of supporting special libraries and archives like
\enquote{The Black Archives} in the Netherlands or the
\enquote{Transgender Archives} in Canada. But the easiest and most
obvious one would be reading banned books.

\emph{MH:} For me, it's important to show people that what's happening
in America isn't just a distant problem. They're ahead of us, but we're
not far behind. I want to make people aware of this. At the same time, I
want to inspire people to take action. We're part of this ourselves and
can take action ourselves. If we don't, we too run the risk of slipping
down that path. The arts, including books, are the first to suffer. We
must take this warning very seriously.

%autor
\begin{center}\rule{0.5\linewidth}{0.5pt}\end{center}

\textbf{Norma Fötsch} is a
librarian at Radboud University in Nijmegen, the Netherlands. In her
spare time, she is working on her PhD thesis. Access to information is a
central theme in her work and research. She is also active in a network
for queer librarians.\newline
\url{https://orcid.org/0009-0009-7009-5520}

\textbf{Marjolein Hordijk} is program director of the Vrijheidscolleges
Foundation, director of the Novio Magica Fantasy Book- and Film
Festival, and a freelancer. She previously worked as a librarian at the
Gelderland Zuid Public Library in the Netherlands. She now remains
connected to the library world through her freelance work.\newline
\url{https://truebeliever-events.nl/}

\textbf{Ivana Smudja} is a film director based in Nijmegen, the Netherlands. She
graduated in 2009 with a degree in film direction from the European
Institute of Design in Milan. Since 2012, she has worked in the
commercial film industry, first as an assistant director and later as a
director. During her formative years, she also worked at Teatro Piccolo
di Milano and Vice Italy, where she developed a documentary-style
approach, blending it with her distinctive artistic vision. Today, Ivana
focuses on socio-historical documentaries that explore topical issues
such as freedom and democracy.\newline
\url{https://www.ivanczyew.com/}

\end{document}